% ---------------------------------------------------------------------
% Immunology sub-section for 05_results_kang_atlas.tex
% Contributed by Immunology agent (agent-mozus3vv) for task t-mozv8ou4
% Target: ~500 words on why Kang 2024 pan-cancer atlas is the right
% stress-test for rare-immune-motif detection. Drop-in ready.
% ---------------------------------------------------------------------

\subsection{Why the Kang pan-cancer immune atlas is the appropriate
stress-test for rare-immune-motif detection}
\label{sec:kang-stress-test-rationale}

The central claim of this work is that our framework detects and
protects rare cell subpopulations that are unannotated before
integration. A credible test of this claim requires a dataset that is
large enough, heterogeneous enough, and clinically consequential
enough that (i) rare subpopulations exist below the detectability
floor of existing methods, (ii) the consequences of their loss are
not academic, and (iii) no single prior annotation can be used as
ground truth, because different source studies have annotated the
same cell types under incompatible lexicons. The pan-cancer immune
atlas of \citet{kang2024integrated} meets all three criteria more
completely than any alternative currently available to the
community.

\textbf{Scale.} With 4.9~million cells across 104 publications, 30
cancer types, and 943 patients, the Kang atlas is approximately
$2\times$ larger than the Human Lung Cell Atlas
\citep{sikkema2023integrated} and an order of magnitude larger than
the pan-cancer myeloid \citep{cheng2021pancancer} or T-cell
\citep{zheng2021pancancer} atlases used in prior integration
benchmarks. At this scale, a rare immune subpopulation present at
$10^{-4}$ prevalence is represented by roughly 500 cells, which is
below the typical cluster-resolution sensitivity of Harmony and
Seurat at default settings but above the information-theoretic
detection floor we derive in \Cref{sec:minimax-detectability}. This
is precisely the regime in which our framework must produce a
decision.

\textbf{Heterogeneity.} The atlas spans 30 cancer types, multiple
tissue microenvironments (primary tumor, metastasis, peripheral
blood, lymph node, adjacent normal), eight sequencing chemistries
(10x Genomics 3' v2, 3' v3, 5' v1, 5' v2; Drop-seq; inDrop; Smart-seq2;
Microwell-seq), and a 943-patient donor covariate that itself
encodes age, sex, therapy line, and geographical origin. Batch
effects in this atlas therefore interact non-trivially: a rare
subpopulation may appear predominantly in 10x 5' v2 fresh-tissue
samples from melanoma patients but be absent from Drop-seq breast
cohorts. This is exactly the cross-batch asymmetry failure mode our
framework is designed to detect (see
\Cref{sec:loss-mode-taxonomy}).

\textbf{Translational stakes.} Rare immune subpopulations in this
atlas are not curiosities. At least five of them --- TPEX
(\textit{CXCR5$^+$TCF7$^+$} stem-like exhausted CD8, predictor of
checkpoint-blockade response \citep{miller2019subsets,
jansen2019expansion}), tumor-reactive CD103$^+$CD39$^+$ tissue-resident
memory CD8 (TIL-therapy selection biomarker
\citep{duhen2018coexpression, caushi2021transcriptional}), effector
regulatory T cells (CCR8-antibody therapy target in Phase~1/2
\citep{desimone2016ttregsig, plitas2016regulatory, vandamme2021depletion}),
LAMP3$^+$ mregDC (LN-drainage marker \citep{maier2020conserved}), and
FOLR2$^+$ perivascular macrophages (breast cancer prognostic
\citep{nalioramos2022tissue}) --- are established clinical biomarkers
or therapy targets whose detection a pan-cancer atlas must support
at scale, not at toy size.

\textbf{Cross-cohort replication as ground truth surrogate.} Because
the atlas aggregates 104 publications, every rare subpopulation we
detect can be independently verified by its replication across
source studies. This supplies a form of cross-cohort ground truth
that does not require prior annotation: a detected rare cluster is
provisionally real if and only if it replicates in $\geq 2$
independent source cohorts with consistent marker signatures
(Jaccard $\geq 0.5$ on the top-20 differentially expressed genes).
This replication-based criterion is what we use to distinguish
genuine previously-unannotated rare subpopulations from integration
artifacts and doublet clusters, following the evidence package
defined in \Cref{sec:evidence-package}.

Taken together, the pan-cancer immune atlas simultaneously provides
the scale at which the detection problem is technically hard, the
heterogeneity that stresses the distribution-shift failure modes
our method targets, the clinical relevance that makes failure
consequential, and the internal replicability that supplies
label-free ground truth. For these reasons we adopt it as the
primary real-world stress-test of our framework, complementing the
controlled label-holdout validations of
\Cref{sec:hlca-ionocyte-holdout,sec:pbmc-asdc-holdout}.

% ---------------------------------------------------------------------
% Bibliography additions required (BibTeX keys used above). See
% workspace/handoffs/bib_additions_immunology.bib for ready-to-paste
% BibTeX entries.
% ---------------------------------------------------------------------
